
% ---------------------------------------------------------------------------- %
%                                  NEW SECTION                                 %
% ---------------------------------------------------------------------------- % 

\chapter{서론}

\section{Excel 인터페이스 개발 목적}

\simulator 은 건물 에너지 시뮬레이션을 위한 강력한 엔진이지만, 그 입력파일 포맷은 컴퓨터가 이해하기 쉬운 방향으로 정리되어 있어, 사용자가 직접 다루기는 다소 불편할 수 있다.
이를 해결하기 위해, \simulator 은 엑셀파일을 입력파일로 사용할 수 있는 인터페이스를 제공한다.
엑셀파일은 많은 사용자에게 친숙한 도구로, 직관적인 데이터 입력과 편집이 가능하여, 비전문가도 쉽게 시뮬레이션 모델을 작성할 수 있다.


\chapter{Excel 인터페이스 사용법}

\section{시트별 관계}
추후 기술

\section{시트별 명세}

\subsection{건물 정보}
추후 기술

\subsection{실}
추후 기술

\subsection{면}
추후 기술

\subsection{개구부}
추후 기술

\subsection{구조체\_면}
추후 기술

\subsection{구조체\_개구부}
추후 기술

\subsection{재료}
추후 기술

\subsection{공급설비}
추후 기술

\subsection{생산설비}
추후 기술

\subsection{환기설비}
추후 기술

\subsection{PV패널}
추후 기술

\chapter{Web GUI 사용법}

\section{실행 및 입력파일 전달}
추후 기술

\section{결과 조회}
추후 기술